\section{Висновки}
Social реалізує вагомий набір взаємодіючих аналітичних компонентів: отримання даних за джерелами, версіоновані знімки, вимірювання інформаційних хвиль, зв'язаний облік доказів, обмежений пошук, типізоване впорядкування подій, локальну обробку медіа та структурні бар'єри звітів. Перевірки на зафіксованих входах відтворюють результати схожості й умовної рандомізації, а явні випробування виявляють семантичні збої, яких великий успішний набір тестів не виключає; звужене правило усуває перевірені збої, не встановлюючи загальної точності випливання. Систему найкраще характеризувати як доказово обґрунтований автоматизований аналітичний фреймворк із нерівномірною валідацією компонентів та обмеженою автономністю процесу. Її практичний внесок --- придатні для перегляду обчислення та збережена невизначеність. Незалежно валідована перевірка фактів, виявлення фізичних подій і загальна наскрізна надійність щодо інцидентів лишаються невстановленими.

\section*{Відтворюваність і дослідницькі артефакти}
Оцінене програмне забезпечення зафіксовано на коміті \code{f0e9d0d67422b0c286a2cb9e7d4e935a833a9533}. Нові файли публікації відокремлені від цієї бази та вказують команди, якими їх створено. Післяаудитна ревізія правила --- коміт \code{319d205}; \code{evaluate_verifier_challenges.py} заново генерує її стрес-тест, а переглянутий пакет із 25 тверджень зберігається поряд із базовим артефактом, не замінюючи його. Супровідний пакет містить аудит архітектури, наукову матрицю «твердження --- докази», знімок середовища, маніфест експериментів, агреговані оцінювальні дані, джерела й входи векторних рисунків, синтетичні артефакти процесу та контрприкладів, а також вихідні файли LaTeX/BibTeX. Точні вхідні значення SHA-256 прив'язують результати до історичного чи розширеного корпусів. Часові тести фіксують зерно 20261106 і 9\,999 збережених вибірок; ця публікація відтворює 25 вибірок на модель і перераховує всі збережені підсумки. Збіг стохастичного префікса відрізняється від свіжої повної симуляції.

Контрактні залежності репозиторію задано діапазонами, а необов'язкові локальні залежності для відео фіксують faster-whisper 1.1.0, PyAV 18.1.0 і requests 2.34.2. Єдиного файла фіксації встановлення немає. Середовище аудиту й середовище формування документів зафіксовано окремо; виміряний запуск тестів використовує Python 3.13.14, а підготовка рисунків і PDF --- системний Python 3.12.3. У чистому клоні коміту \code{38fe93d}, який не містить ігнорованих Git корпусів, пройшли 788 із 794 тестів; чотири збої та дві помилки спричинені цими відсутніми локальними файлами. Повний повторний історичний збір із платформ не оцінювали. Практики відтворюваності дотримуються принципу збереження даних, перетворень, параметрів і проміжних результатів \citep{sandve2013reproducible,wilson2017practices}; приватний корпус унеможливлює необмежене відтворення всіх етапів третіми сторонами. Принципи FAIR обґрунтовують придатні для повторного використання метадані, не вимагаючи, щоб чутливі щодо приватності дані були відкрито завантажуваними \citep{wilkinson2016fair}.

\section*{Доступність даних і коду}
Вихідний код публічно доступний за адресою \url{https://github.com/cyberdid/social-framework} (тег v1.0.0, коміт \code{b48cc97}) за ліцензією Apache License 2.0. Цей випуск містить вихідні модулі, схеми та 694 тести, які працюють без приватних даних, ще шість тестів пропущено; також проходять 16 підтестів. Він не містить дослідницьких корпусів, звітів, кураторських збирачів кейсу, входів стрес-тесту й класифікацій каналів, бо вони містять персональні дані або матеріали кейсу про людей, яких можна ідентифікувати, тож наведені тут результати не можна повністю відтворити лише з нього. Репозиторій розробки лишається приватним. Стаття поширюється за ліцензією CC BY 4.0. Сирі дані Telegram, коментарів і медіа та облікові дані разом зі статтею не розповсюджуються. Агреговані й синтетичні додаткові матеріали, що зберігають приватність, супроводжують рукопис у каталозі публікації й можуть окремо розглядатися для оприлюднення. Окремий рукопис кейсу лишається проєктною чернеткою з непідтвердженим статусом публікації. Доступність коду, обчислювальна повторюваність і незалежна семантична валідація є окремими обмеженнями.

\section*{Внесок автора}
Олександр Пилипенко керував концептуалізацією, проєктуванням системи, дизайном дослідження та напрямом реалізації; спрямовував аналіз та інтерпретацію; керував підготовкою рукопису. Це твердження не означає особистого виконання кожного програмного чи обчислювального кроку. Автор відповідає за остаточну перевірку й затвердження.

\section*{Фінансування та конфлікт інтересів}
Дослідження не отримало зовнішнього фінансування. Автор заявляє про відсутність конфлікту інтересів відповідно до останніх підтверджених декларацій рукописів проєкту.

\section*{Розкриття щодо генеративного ШІ}
Інструменти генеративного ШІ допомагали в розробці програмного забезпечення, документуванні, обчислювальному аудиті, упорядкуванні літератури, підготовці й редагуванні рукопису та його перекладі українською мовою. Структурований огляд методології, архітектури, доказів і редакції виконала одна ШІ-система в межах одного сеансу; це не є незалежним людським рецензуванням, еталонною розміткою чи підтвердним доказом. Історичний зовнішній пілот Gemini на ключових кадрах описано лише в межах його задокументованого обсягу; для наведених обчислень публікаційного аудиту нових платних викликів моделей не використовували. Автор зберігає відповідальність за перевірку, тлумачення та затвердження рукопису і його дослідницьких артефактів.
